% 第一章（1a）附录：本部分正文引用的推导
\section{第一章学术部分的推导}\label{app:ch1a}

\subsection{正态曲线与 Elo 逻辑斯蒂曲线的刻度换算}\label{app:ch1a-scale}

Elo 曲线\cref{eq:1a-elo-curve} 在分差 0 处的斜率为 $\Elogit'(0)=\cElo/4=\ln 10/1600\approx 1.439\times10^{-3}$（每分）。正态曲线 $\Phi\big(x/(\sigma\sqrt2)\big)$ 在 0 处的斜率为 $1/(\sigma\sqrt2\sqrt{2\pi})$。令两者相等，得
\[
  \sigma=\frac{1600}{\ln 10\cdot\sqrt{2}\cdot\sqrt{2\pi}}\approx 196 ,
\]
即“400 分刻度的逻辑斯蒂曲线”在中心附近等价于“每人单局表现标准差约 196 分”的正态曲线，与 Elo 原书取的 200 几乎相同\tagOwn。用 FIDE 现行条例 8.1.1 表（$p=0.51,\dots,0.99$ 共 49 格）做最小二乘拟合，得到的单人标准差为 202.6；表中分差与 $\sigma=200$ 的正态曲线最大相差 19 分（$p=0.99$ 处），与 Elo 逻辑斯蒂曲线最大相差 121 分（同一处）\tagOwn。两种曲线在中段几乎重合，差别只在 $p>0.9$ 的尾部：逻辑斯蒂曲线给弱者留的胜率更大。

\subsection{TrueSkill 两人更新式}\label{app:ch1a-ts}

设胜方、负方实力的先验为 $s_{\mathrm w}\sim\Normal(\mu_{\mathrm w},\sigma_{\mathrm w}^2)$、$s_{\mathrm l}\sim\Normal(\mu_{\mathrm l},\sigma_{\mathrm l}^2)$，表现 $p=s+\Normal(0,\beta^2)$。表现差 $d=p_{\mathrm w}-p_{\mathrm l}\sim\Normal(m,c^2)$，其中 $m=\mu_{\mathrm w}-\mu_{\mathrm l}$，$c^2=\sigma_{\mathrm w}^2+\sigma_{\mathrm l}^2+2\beta^2$。观测到“胜方获胜且不是平局”即 $d>\varepsilon$。

\paragraph{截断正态的矩。}记 $d=m+cZ$，$Z$ 为标准正态，$t=m/c$，$a=\varepsilon/c$，则 $d>\varepsilon\Leftrightarrow Z>-(t-a)$。对 $u=t-a$，
\[
  \E[Z\mid Z>-u]=\frac{\varphi(u)}{\Phi(u)}=:v(t,a),\qquad
  \Var[Z\mid Z>-u]=1-v(t,a)\big(v(t,a)+u\big)=:1-w(t,a).
\]
\paragraph{传回实力。}$(s_{\mathrm w},d)$ 联合正态，$\Cov(s_{\mathrm w},d)=\sigma_{\mathrm w}^2$，故 $\E[s_{\mathrm w}\mid d]=\mu_{\mathrm w}+(\sigma_{\mathrm w}^2/c^2)(d-m)$，$\Var[s_{\mathrm w}\mid d]=\sigma_{\mathrm w}^2(1-\sigma_{\mathrm w}^2/c^2)$。对条件 $d>\varepsilon$ 取期望（全期望与全方差公式）：
\[
  \E[s_{\mathrm w}\mid d>\varepsilon]=\mu_{\mathrm w}+\frac{\sigma_{\mathrm w}^2}{c}\,v(t,a),\qquad
  \Var[s_{\mathrm w}\mid d>\varepsilon]=\sigma_{\mathrm w}^2\Big(1-\frac{\sigma_{\mathrm w}^2}{c^2}\Big)+\frac{\sigma_{\mathrm w}^4}{c^4}\,c^2\big(1-w\big)
  =\sigma_{\mathrm w}^2\Big(1-\frac{\sigma_{\mathrm w}^2}{c^2}\,w(t,a)\Big).
\]
负方的 $\Cov(s_{\mathrm l},d)=-\sigma_{\mathrm l}^2$，均值更新取负号，方差更新形式相同。TrueSkill 把后验近似为具有这两个矩的高斯（期望传播中的矩匹配），即得\cref{eq:1a-ts-update}。若有动态方差 $\gamma^2$（TrueSkill 论文的记号），赛前先令 $\sigma^2\leftarrow\sigma^2+\gamma^2$。爆冷时 $t\ll0$，$v\approx a-t$ 很大，均值更新幅度近似与“爆冷程度”成正比；$t\gg0$（强者赢弱者）时 $v\to0$，几乎不更新\tagOwn。

\subsection{TrueSkill 匹配质量即相对平局概率}\label{app:ch1a-quality}

平局边际 $\varepsilon$ 很小时，$\Prob(|d|\le\varepsilon)\approx 2\varepsilon\cdot\frac{1}{c\sqrt{2\pi}}\exp\!\big(-\frac{m^2}{2c^2}\big)$。它的最大可能值出现在 $m=0$ 且双方实力完全确定（$\sigma_i=\sigma_j=0$，$c^2=2\beta^2$）时，为 $2\varepsilon/\big(\sqrt{2\beta^2}\sqrt{2\pi}\big)$。两者之比与 $\varepsilon$ 无关：
\[
  q_{\text{draw}}=\sqrt{\frac{2\beta^2}{c^2}}\exp\!\Big(-\frac{m^2}{2c^2}\Big),\qquad c^2=2\beta^2+\sigma_i^2+\sigma_j^2 ,
\]
即\cref{eq:1a-ts-quality}。它同时惩罚均值差（指数项）和不确定性（根号项）：两个新号即使均值相同，质量也低于两个均值相同的老玩家\tagOwn。

\subsection{EOMM 的三种特例}\label{app:ch1a-eomm}

设玩家流失只取决于本局结果，记 $c(\mathrm W)$、$c(\mathrm L)$、$c(\mathrm D)$ 为胜、负、平后的流失概率，$P_{ij}(\cdot)$ 为 $i$ 对 $j$ 的结果分布。一对玩家的流失之和为
\[
  P_{ij}(\mathrm W)\big[c(\mathrm W)+c(\mathrm L)\big]+P_{ij}(\mathrm L)\big[c(\mathrm L)+c(\mathrm W)\big]+P_{ij}(\mathrm D)\cdot 2c(\mathrm D)
  =\big[c(\mathrm W)+c(\mathrm L)\big]-P_{ij}(\mathrm D)\big[c(\mathrm W)+c(\mathrm L)-2c(\mathrm D)\big].
\]
对所有配对求和后，最小化总流失等价于：若 $c(\mathrm W)+c(\mathrm L)>2c(\mathrm D)$，最大化总平局概率；若不等号反向，最小化总平局概率。平局在实力相近时最可能发生，所以前一种情形下按分匹配最优，后一种情形下按分匹配最差。若流失与本局结果无关（只取决于赛前状态），上式与配对无关，任何匹配（包括随机匹配）都一样。这就是 Chen 等\cite{chen2017eomm} 第 4 节的结论。

由此还可以看出一点（\tagOwn 的推论）：在\emph{没有平局}的输赢游戏中，若流失只取决于本局结果，则每对玩家的流失之和恒为 $c(\mathrm W)+c(\mathrm L)$，匹配对总流失毫无影响。参与度优化匹配要有效果，流失必须依赖于玩家的状态（例如最近几局的胜负、技能、在线时长），而且这种依赖要因人而异。EOMM 的输入正是“最近 10 局的胜负平”等状态变量，这也是它与“按连胜连败安排对手”在机制上相通的原因。

\subsection{Yuksel 目标胜率公式}\label{app:ch1a-yuksel}

设玩家最近 $n$ 局赢了 $w$ 局，希望包含接下来 $n+1$ 局在内、长 $2n+1$ 的窗口的期望胜率为 $\lambda$。若接下来每局对手都使其胜率为 $p$，窗口期望胜局数为 $w+(n+1)p$，令其等于 $\lambda(2n+1)$ 即得\cref{eq:1a-yuksel}。两人游戏取 $\lambda=1/2$、论文实验取 $n=5$ 时，$w=0,1,\dots,5$ 对应的目标胜率依次为 $11/12,\,9/12,\,7/12,\,5/12,\,3/12,\,1/12$（约 92\%、75\%、58\%、42\%、25\%、8\%）\tagOwn。论文另取概率窗口 $\Delta p=1/(2n+2)=1/12$，在评分 $f^{-1}(p\pm\Delta p)$ 范围内找对手；这保证 $p-\Delta p>0$ 且 $p+\Delta p<1$\cite{yuksel2024skill}。换言之，这一匹配器对近期战绩的反应相当强：连赢 5 局后，下一局几乎一定会遇到明显更强的对手。
